% Evans 10.2: printed pp.550--554; English PDF pp.567--571.
% Chinese section start: printed p.389 / scan p.407.
\Needspace{80mm}
\section{Uniqueness}
\label{sec:10-2}
\subsection{Tests at the upper time boundary}
\Needspace{50mm}
\begin{lemma}[Terminal extrema for an initial-value solution;
source \S10.2 lemma]
Let $u$ be a continuous viscosity subsolution on $0<t<T$.
If $u-\phi$ has a local maximum at $(x_0,T)$ relative to $t\leq T$,
then $\phi_t+H(D\phi,x_0)\leq0$ there. A supersolution has the
reverse inequality at a relative local minimum.
\end{lemma}
\begin{expandedproof}[discard a positive penalty derivative]
Strictify the maximum and restrict to a small closed backward cylinder.
For $t<T$ set $\phi_\delta=\phi+\delta/(T-t)$.
The quantity $u-\phi_\delta$ tends to $-\infty$ as $t\uparrow T$.
A comparison at $(x_0,T-\sqrt\delta)$ shows its maximum tends to
$(u-\phi)(x_0,T)$. The strict gap on the other cylinder boundaries
therefore puts maximizing points $(x_\delta,t_\delta)$ in the open
PDE domain, and all cluster points equal $(x_0,T)$. At those points,
\[
 \phi_t+{\delta\over(T-t_\delta)^2}+H(D\phi,x_\delta)\leq0.
\]
Discard the nonnegative middle term and pass to the limit. It need not
tend to zero. For a lower contact use
$\phi_\delta=\phi-\delta/(T-t)$, and discard its nonpositive derivative
from the reverse inequality. This lemma does not reverse the viscosity
convention: it only extends the initial-value convention to a one-sided
contact at the top of its slab.
\end{expandedproof}
\Needspace{42mm}
\PDEAnnotation{终端接触与终值问题不同}{本引理研究初值解在 $t=T$ 的单侧接触，不改变不等式方向。下一节的终值控制问题采用反向时间演化，才会改变方向；这两个“终端”不能混用。}

\Needspace{70mm}
\subsection{Comparison by doubling space and time}
Assume $H$ satisfies, with a single constant $C$,
\begin{equation}\label{eq:10-H-lip}
 \begin{split}
 |H(p,x)-H(q,x)|&\leq C|p-q|,\\
 |H(p,x)-H(p,y)|&\leq C|x-y|(1+|p|).
 \end{split}
\end{equation}
\Needspace{55mm}
\begin{theorem}[Comparison; source Theorem 1, strengthened to sub/supersolutions]
Let $u,v\in\mathrm{BUC}(\R^n\times[0,T])$ be a subsolution and a
supersolution respectively, with the initial-value convention.
If $u(x,0)\leq v(x,0)$, then $u\leq v$ on the slab. In particular
there is at most one BUC viscosity solution with prescribed initial data.
More generally,
\[
 \sup_x(u(x,t)-v(x,t))\leq\sup_x(u(x,0)-v(x,0)).
\]
\end{theorem}
\begin{expandedproof}[localization, small penalty, and two touching tests]
\textit{1. A positive maximum after penalization.}
Suppose $u(x_*,t_*)-v(x_*,t_*)=a>0$. Fix $\eta>0$ with
$2\eta T<a/4$. For $\varepsilon>0$ maximize on
$\R^{2n}\times[0,T]^2$
\begin{equation}\label{eq:10-doubling}
 \begin{split}
 \Phi_\varepsilon(x,y,t,s)={}&u(x,t)-v(y,s)-\eta(t+s)\\
 &-{|x-y|^2+|t-s|^2\over2\varepsilon}
 -\varepsilon(|x|^2+|y|^2).
 \end{split}
\end{equation}
Boundedness and the quadratic confinement make the maximum attained;
call its point $(x_\varepsilon,y_\varepsilon,t_\varepsilon,s_\varepsilon)$.
For small $\varepsilon$, testing on $(x_*,x_*,t_*,t_*)$ shows that
this maximum is at least $a/2$. Below all variables refer to this point.

\textit{2. First bounds.}
Comparing with a fixed point, and using boundedness of $u,v$, gives
\[
 r:=|x-y|=O(\sqrt\varepsilon),\quad
 d:=|t-s|=O(\sqrt\varepsilon),\quad
 \varepsilon(|x|+|y|)=O(\sqrt\varepsilon).
\]
These bounds allow $x,y$ to escape to infinity. This is why uniform
continuity on the full slab, rather than just on compact sets, is used.

\textit{3. The stronger small-penalty estimate.}
Compare the maximum with $\Phi_\varepsilon(x,x,t,t)$.
If $\omega_v$ is a modulus for $v$, then
\[
 {r^2+d^2\over2\varepsilon}
 \leq\omega_v(r+d)+\eta d
       +\varepsilon(|x|+|y|)r=o(1).
\]
In particular $r^2/\varepsilon\to0$. Merely knowing
$r=O(\sqrt\varepsilon)$ would only bound this expression and would
not complete the proof.

\textit{4. Neither time is initial.}
Since the penalty is nonnegative, $u(x,t)-v(y,s)\geq a/2$.
If a subsequence had $t\to0$, then $s\to0$ as well, and initial
ordering plus uniform continuity would give
\[
 u(x,t)-v(y,s)\leq\omega_u(t)+\omega_v(r+s)\longrightarrow0.
\]
This is a contradiction. The same reasoning excludes $s\to0$.
Thus both contact times are positive for all small $\varepsilon$.
They may equal $T$, in which case the preceding lemma applies.

\textit{5. Freeze the other pair of variables.}
At the maximum, the penalty with $(y,s)$ frozen is an upper test for
$u$ after adding a constant. Its time derivative and spatial gradient are
\[
 b_u=\eta+{t-s\over\varepsilon},\qquad
 p_u={x-y\over\varepsilon}+2\varepsilon x.
\]
The negative penalty with $(x,t)$ frozen is a lower test for $v$, with
\[
 b_v=-\eta+{t-s\over\varepsilon},\qquad
 p_v={x-y\over\varepsilon}-2\varepsilon y.
\]
Consequently $b_u+H(p_u,x)\leq0$ and $b_v+H(p_v,y)\geq0$.
Subtracting, the uncontrolled time difference cancels:
\[
 2\eta\leq H(p_v,y)-H(p_u,x).
\]

\textit{6. Estimate the Hamiltonians.}
By \eqref{eq:10-H-lip},
\[
 \begin{split}
 2\eta&\leq C|p_u-p_v|+Cr(1+|p_v|)\\
 &\leq C\sqrt\varepsilon+Cr
        +C{r^2\over\varepsilon}+Cr\sqrt\varepsilon\longrightarrow0.
 \end{split}
\]
This contradicts the fixed $\eta>0$, proving $u\leq v$.
Adding a constant to $v$ leaves its supersolution property unchanged;
apply ordering to $v+\sup_x(u(x,0)-v(x,0))$ for the last assertion.
Applying the result in both directions yields uniqueness and sup-norm
contraction for solutions with different initial values.
\end{expandedproof}
\PDEBodyDiagram{10.2}{Doubling variables produces two independent smooth tests.}
{The common vector $(x-y)/\varepsilon$ need not be bounded. Global Lipschitz control in momentum handles the small difference $p_u-p_v$, while $|x-y|^2/\varepsilon\to0$ handles the spatial dependence of $H$. The opposite strict time shifts leave the positive gap $2\eta$ after subtraction.}
{ch10-doubling}
\Needspace{42mm}
\PDEAnnotation{小量必须足够强}{若只得到 $|x-y|=O(\sqrt\varepsilon)$，则 $|x-y|^2/\varepsilon$ 可能不趋零。对角比较给出整个罚项 $o(1)$，正好消除 Hamiltonian 的空间误差。}

\Needspace{70mm}
\subsection{A comparison extension for Hopf--Lax}
The quadratic Hamiltonian $H(p)=|p|^2/2$ is not globally Lipschitz.
The source's Hopf--Lax uniqueness statement therefore needs a separate
comparison observation, supplied here.
\Needspace{55mm}
\begin{theorem}[One spatially Lipschitz competitor; editorial supplement]
Suppose $H(p)$ is continuous and independent of $x$. The comparison
conclusion above still holds if at least one of $u,v$ is uniformly
Lipschitz in $x$ on the slab. The other function may be merely BUC.
\end{theorem}
\begin{expandedproof}[a contact bounds the momentum]
If a $K$-Lipschitz function has a differentiable upper or lower touching
test at a point, its spatial test gradient has norm at most $K$.
Indeed restrict to increments $he$, use the touching inequality and the
Lipschitz bound, divide by $h$ with both signs, and conclude
$|D\phi\cdot e|\leq K$ for all unit $e$.

Use the same doubled maximum and first four steps above. If $u$ is
Lipschitz then $|p_u|\leq K$; if $v$ is Lipschitz then $|p_v|\leq K$.
Since $|p_u-p_v|=O(\sqrt\varepsilon)$, both momenta lie in a fixed
compact ball. Uniform continuity of $H$ on that ball gives
$H(p_v)-H(p_u)\to0$, contradicting $2\eta>0$ exactly as before.
No global Lipschitz bound on $H$, convexity, or boundedness of the
other function's gradient is needed for this particular extension.
\end{expandedproof}
\Needspace{42mm}
\PDEAnnotation{超线性 Hamiltonian 的衔接}{$H(p)=|p|^2/2$ 的梯度无界，不能直接套用全局动量 Lipschitz 定理。若一方解在空间中为 $K$-Lipschitz，接触斜率被限制在有界球内，就只需该球上的一致连续性。}
